\chapter{Rash Evaluation}
\section*{What Is This Test?} Rash evaluation identifies the morphology, distribution, timing, and possible cause of a skin eruption.\section*{Alternative Names} Skin-rash assessment; dermatologic examination; eruption evaluation.\section*{Principle} History and inspection guide targeted testing such as swabs, scrapings, blood tests, biopsy, or allergy assessment.\section*{Why This Test Is Ordered} It evaluates new, spreading, painful, itchy, blistering, infected, medication-related, or recurrent rashes.\section*{Primary vs Secondary Test} Examination is primary; tests are selected for the suspected cause. Fever, mucosal lesions, facial swelling, or skin pain can signal urgency.\section*{How This Test Is Performed} The clinician examines skin, nails, hair, mouth, and lymph nodes and documents distribution and morphology.\section*{What Preparation Is Needed} Avoid applying new creams before evaluation and bring medicine, exposure, and photograph history.\section*{Understanding Results} Many rashes overlap in appearance; a normal screening test does not exclude all skin disease.\section*{Follow-up Tests} KOH scraping, bacterial or viral swab, CBC, serology, patch testing, biopsy, or dermatology referral may follow.\section*{References} \begin{enumerate}\item MedlinePlus. Rash Evaluation.\item American Academy of Dermatology. Rash evaluation resources.\end{enumerate}\clearpage\section*{Understanding Results and Follow-up}
The laboratory report should identify the specimen, method, units, reference interval or decision limit, and whether the result is positive, negative, abnormal, or inconclusive. Reference intervals vary with age, sex, pregnancy, specimen type, assay, and laboratory; a result outside a range is not automatically a diagnosis, and a result within a range does not always exclude disease. Discuss unexpected results with the ordering clinician, who can decide whether repeat or confirmatory testing, treatment, or referral is appropriate. Do not change medicines or delay urgent care based on an isolated result.
\section*{Regulatory Status and Authorization Date}This chapter describes a test category, not a specific commercial product. FDA, CDSCO, EU IVDR, and MHRA approval or registration dates therefore cannot be assigned without the assay manufacturer, product name, instrument, and intended use. For laboratory-developed tests, an individual agency approval date may not exist. See the Preface for the interpretation of regulatory status.
\clearpage
